\documentclass[conference]{IEEEtran}

\usepackage{amsmath,amssymb,amsfonts}
\usepackage{textcomp}
\usepackage{xcolor}
\usepackage{booktabs}
\usepackage{hyperref}
\usepackage{url}

\hypersetup{colorlinks=true, linkcolor=blue, citecolor=blue, urlcolor=blue}

\begin{document}

\title{Midpoint Report: Benchmarking KV Cache Quantization and Token-Level Selection for Efficient LLM Inference}

\author{
\IEEEauthorblockN{Suhas Morisetty\textsuperscript{*}, Rishika Mamidibathula\textsuperscript{\dag}, Ziheng Wang\textsuperscript{*}, Tony Tian\textsuperscript{*}}
\IEEEauthorblockA{
\textsuperscript{*}Department of Computer Science, Columbia University \\
\textsuperscript{\dag}Department of Data Science, Columbia University \\
COMS E6998: High Performance Machine Learning \\
\{vm2825, rm4318, zw3153, jt3640\}@columbia.edu}
}

\maketitle

\begin{abstract}
Autoregressive LLM inference is bottlenecked by the key-value (KV) cache, whose memory grows linearly with sequence length. Published optimizations use inconsistent evaluation setups, making cross-method comparison unreliable. We benchmark two representative techniques---KIVI asymmetric quantization and TopK (TokenSelect) dynamic token selection---against an FP16 baseline on LLaMA-2-7B, with shared weights, prompts, seeds, and hardware (A100-80GB). Across 72 synthetic runs at three context lengths, 720 LongBench evaluations, and 6 perplexity measurements, we find that KIVI 4-bit achieves 2.90$\times$ compression with no quality loss (PPL 6.93 vs.\ 6.94), while 2-bit compresses 4.47$\times$ at nearly doubled perplexity (13.70). TopK preserves quality exactly but provides no memory reduction, as it sparsifies attention over a full cache. Both methods incur throughput overhead in pure PyTorch that would diminish with the original authors' CUDA kernels.
\end{abstract}

\begin{IEEEkeywords}
KV cache, large language models, inference optimization, quantization, token selection
\end{IEEEkeywords}

%% ============================================================
\section{Project Summary}
%% ============================================================

\textbf{Title.} Benchmarking KV Cache Quantization and Token-Level Selection for Efficient LLM Inference.

\textbf{Goals.} Fair, controlled comparison of KV cache optimization techniques under identical conditions. We implement methods from two optimization families---quantization and dynamic token selection---within the same inference framework, so observed differences are attributable solely to the cache policy. All methods operate at inference time with no retraining or architectural changes. We evaluate memory, throughput, latency, perplexity (WikiText-103), and downstream task performance (LongBench F1/ROUGE-L).

\textbf{Model.} LLaMA-2-7B~\cite{llama2} in FP16: 32 layers, 32 heads, head dim 128. KV cache at 8192 tokens $\approx$ 4.4~GB.

\textbf{Datasets.}
\begin{itemize}
    \item \textit{Synthetic prompts}: 12 WikiText-103~\cite{wikitext} passages (4 per target length: 512, 2048, 8192 tokens) wrapped with a summarization instruction. Used for throughput, latency, memory, and compression measurements.
    \item \textit{WikiText-103 test}: 20 passages for perplexity via a split-text protocol---all methods prefill the first half, apply their cache modification, and are scored on the second half under identical conditions.
    \item \textit{LongBench}~\cite{longbench}: 6 long-context tasks $\times$ 20 examples each (QA scored with token-level F1, summarization with ROUGE-L).
\end{itemize}

\subsection{Performance Optimization Techniques}

\textbf{Baseline (FP16 Full Cache).} Standard dense attention over all past tokens with no cache modification. Serves as the reference for all comparisons.

\textbf{KIVI: Asymmetric Block-Wise Quantization.} KIVI~\cite{kivi} reduces KV cache precision from FP16 to 4-bit or 2-bit. Keys are quantized per-channel and values per-token, reflecting their distinct outlier distributions (keys have channel-axis outliers, values have token-axis outliers).

Our implementation uses block-wise group quantization with group\_size=32. Rather than re-quantizing the entire history when new tokens arrive (O($N$) per step), we partition the quantized cache into independent blocks, each with its own scale and zero-point. New tokens accumulate in an FP16 overflow buffer; once it reaches group\_size, the buffer is quantized as a new block. Existing blocks are never re-quantized, giving O(1) amortized cost. A sliding window of 128 recent tokens remains in FP16 to preserve short-range attention fidelity.

\textbf{TopK: Dynamic Token-Level Selection (TokenSelect).} TopK~\cite{topk} maintains the full FP16 cache but restricts each decode step's attention to a dynamically chosen subset. Our implementation includes the key contributions from the TokenSelect paper:
\begin{itemize}
    \item \textit{Q-dependent scoring}: tokens scored via $q \cdot K^\top$ at each step, making selection dependent on the current query.
    \item \textit{Head-level soft voting}: scores aggregated across heads with per-head criticality weights for consensus selection.
    \item \textit{Paged KV storage}: cache organized in 64-token pages for efficient non-contiguous retrieval.
    \item \textit{Selection cache reuse}: previous indices reused when consecutive queries have $>$0.95 cosine similarity.
    \item \textit{Sink + local tokens}: 128 initial tokens and 512 recent tokens always included.
    \item \textit{Periodic full-context refresh}: full cache attention every \texttt{refresh\_interval} steps to prevent drift.
\end{itemize}

\textbf{Unified design.} Both methods conform to a shared \texttt{MethodWrapper} interface. The generation loop is identical across all methods; only the cache transformation differs. All implementations use pure PyTorch without custom CUDA/Triton kernels, isolating algorithmic overhead from kernel engineering.

\subsection{Profiling Tools}
GPU memory via \texttt{torch.cuda.max\_memory\_allocated()}; timing via \texttt{torch.cuda.synchronize()} + \texttt{time.perf\_counter()} for TTFT and per-token latency; per-method \texttt{get\_kv\_size\_bytes()} for true cache footprint at actual bit-widths; Weights \& Biases for logging; Modal A100-80GB containers for compute.

%% ============================================================
\section{Current Progress}
%% ============================================================

\textbf{Links.} GitHub: \url{https://github.com/rishika1099/KV-Cache-Implementation} \\ W\&B: \url{https://wandb.ai/rm4318-columbia-university/kv-cache-benchmark/runs/rtho04qc?nw=nwuserrm4318}

The full benchmark pipeline is operational end-to-end: prompt generation, method dispatch, metric collection, LongBench scoring, WikiText perplexity evaluation, result merging, and automated plot generation. We completed 72 synthetic inference runs (6 configs $\times$ 12 prompts), 720 LongBench evaluations (6 configs $\times$ 6 tasks $\times$ 20 examples), and 6 perplexity measurements on Modal A100-80GB GPUs.

\subsection{Main Results}

Table~\ref{tab:main_results} summarizes aggregate metrics. KIVI 4-bit achieves an average compression ratio of 2.90$\times$ (mean KV cache 626.4~MB vs.\ baseline 1990.7~MB) with essentially no quality loss (PPL 6.93 $\approx$ baseline 6.94). The 4-bit quantization grid's 16 discrete levels introduce rounding error below the model's noise floor, and block-wise quantization (independent scale/zero per 32-token block) helps by adapting to local value distributions. KIVI 2-bit achieves stronger 4.47$\times$ compression but nearly doubles perplexity to 13.70---only 4 discrete levels cause meaningful distortion of key-value representations.

TopK reports 1.00$\times$ compression across all K values, consistent with its design of maintaining the complete cache while sparsifying only the attention computation. All TopK configurations match baseline perplexity exactly (6.94), since the full cache is maintained and the selection mechanism always includes sink tokens, local tokens, and the highest-scoring positions.

\begin{table}[t]
\centering
\caption{Aggregate results averaged across all sequence lengths and prompts.}
\label{tab:main_results}
\begin{tabular}{@{}llrrrr@{}}
\toprule
\textbf{Method} & \textbf{Config} & \textbf{Comp.} & \textbf{TPS} & \textbf{Lat.(ms)} & \textbf{PPL} \\
\midrule
Baseline & FP16 & 1.00$\times$ & 36.1 & 27.9 & 6.94 \\
\midrule
KIVI & 4-bit, grp=32 & 2.90$\times$ & 8.7 & 220.5 & 6.93 \\
KIVI & 2-bit, grp=32 & 4.47$\times$ & 8.4 & 235.7 & 13.70 \\
\midrule
TopK & K=512 & 1.00$\times$ & 27.9 & 37.5 & 6.94 \\
TopK & K=1024 & 1.00$\times$ & 27.3 & 38.4 & 6.94 \\
TopK & K=2048 & 1.00$\times$ & 32.3 & 33.5 & 6.94 \\
\bottomrule
\end{tabular}
\end{table}

\subsection{Compression Scaling with Sequence Length}

Table~\ref{tab:seq_scaling} shows KIVI's compression improving at longer contexts. The fixed 128-token FP16 residual window constitutes $\sim$25\% of the cache at 512 tokens but under 2\% at 8192, so the quantized fraction---and therefore the compression ratio---grows with sequence length. At 8192 tokens, 4-bit reaches 3.33$\times$ and 2-bit reaches 5.63$\times$. TopK remains at 1.00$\times$ at all lengths.

\begin{table}[t]
\centering
\caption{KV cache size (MB) and compression ratio by sequence length.}
\label{tab:seq_scaling}
\begin{tabular}{@{}lrrrrrr@{}}
\toprule
 & \multicolumn{2}{c}{\textbf{512 tok}} & \multicolumn{2}{c}{\textbf{2048 tok}} & \multicolumn{2}{c}{\textbf{8192 tok}} \\
\cmidrule(lr){2-3} \cmidrule(lr){4-5} \cmidrule(lr){6-7}
\textbf{Method} & \textbf{MB} & \textbf{Comp.} & \textbf{MB} & \textbf{Comp.} & \textbf{MB} & \textbf{Comp.} \\
\midrule
Baseline & 380 & 1.00$\times$ & 1185 & 1.00$\times$ & 4407 & 1.00$\times$ \\
KIVI 4b & 161 & 2.36$\times$ & 394 & 3.01$\times$ & 1325 & 3.33$\times$ \\
KIVI 2b & 122 & 3.11$\times$ & 254 & 4.66$\times$ & 783 & 5.63$\times$ \\
TopK & 380 & 1.00$\times$ & 1185 & 1.00$\times$ & 4407 & 1.00$\times$ \\
\bottomrule
\end{tabular}
\end{table}

\subsection{Throughput and Latency}

TopK at K=2048 reaches 32.3 TPS (89\% of baseline's 36.1 TPS). Smaller K values are counterintuitively slower (27--28 TPS): the token selection overhead (scoring all positions, head voting, paged retrieval) dominates the compute savings from attending to fewer tokens at our tested sequence lengths. At longer sequences where full attention cost dominates, smaller K would show greater relative speedup.

KIVI runs at 8.4--8.7 TPS ($\sim$76\% slower) due to the per-step dequantize-attend cycle: all quantized blocks are expanded to FP16, concatenated with the residual window, used for attention, and the new token routed to the overflow buffer. In pure PyTorch this involves multiple memory allocations and element-wise arithmetic. The original KIVI paper~\cite{kivi} achieves near-baseline throughput with fused Triton kernels that perform quantized attention without full dequantization. TTFT follows the same pattern: baseline 460~ms, TopK 323--487~ms, KIVI 1.6s.

\subsection{LongBench Downstream Tasks}

Table~\ref{tab:longbench} presents per-task scores. KIVI 4-bit (avg 0.062) slightly exceeds baseline (0.060), confirming that 4-bit quantization does not degrade downstream task quality. TopK at K$\geq$1024 matches baseline exactly; K=512 drops to 0.052, suggesting that restricting attention to only 512 tokens occasionally causes the model to miss relevant mid-document information.

KIVI 2-bit shows the most variation: it performs comparably on QA tasks (short, localized answers) but degrades on summarization (gov\_report, qmsum) where the model must synthesize across the full document---consistent with its elevated perplexity. Overall scores are modest across all methods including baseline, reflecting the inherent difficulty of these tasks for a 7B model with greedy decoding and 200-token generation limit.

\begin{table}[t]
\centering
\caption{LongBench scores (F1 for QA, ROUGE-L for summarization), 20 examples per task.}
\label{tab:longbench}
\begin{tabular}{@{}lrrrrr@{}}
\toprule
\textbf{Task} & \textbf{Base} & \textbf{KIVI-4} & \textbf{KIVI-2} & \textbf{TopK-512} & \textbf{TopK-1024} \\
\midrule
qasper & .051 & .055 & .044 & .051 & .051 \\
multifieldqa & .049 & .050 & .050 & .043 & .049 \\
hotpotqa & .012 & .011 & .015 & .013 & .012 \\
2wikimqa & .020 & .023 & .043 & .018 & .020 \\
gov\_report & .154 & .161 & .116 & .135 & .154 \\
qmsum & .074 & .069 & .036 & .053 & .074 \\
\midrule
\textbf{Average} & \textbf{.060} & \textbf{.062} & \textbf{.050} & \textbf{.052} & \textbf{.060} \\
\bottomrule
\end{tabular}
\end{table}

\subsection{Peak GPU Memory}

An important nuance is that KIVI's peak GPU memory (24.9~GB average) exceeds baseline (20.5~GB) despite achieving the highest cache compression. This occurs because the per-step dequantization temporarily creates full FP16 tensors that coexist with the quantized cache, overflow buffer, and residual window. The peak reflects these transient allocations, not the steady-state compressed footprint. With fused kernels operating directly on quantized representations, this overhead would be eliminated. TopK's peak memory is close to baseline at K=512 (21.6~GB) and K=1024 (22.0~GB), but K=2048 shows elevated peak memory (35.3~GB) due to larger intermediate tensors in the attention computation.

\subsection{Key Observations}

The results reveal that KIVI and TopK occupy complementary positions in the tradeoff space. KIVI reduces the physical size of the stored cache, delivering genuine memory savings that improve with sequence length---at 8192 tokens, 4-bit KIVI stores the KV cache in 1.3~GB versus baseline's 4.4~GB. However, it pays for this with throughput overhead from the dequantize-attend cycle, and at aggressive bit-widths (2-bit), with quality degradation. TopK takes the opposite approach: it stores everything but computes less, preserving quality exactly while modestly reducing per-step compute. Neither method dominates on all axes, confirming that the choice of KV cache policy depends on whether the deployment prioritizes memory footprint, throughput, or output fidelity.

A second observation is that block-wise group quantization is critical to KIVI's practical viability. Without it, every new token forces a full re-quantization of the entire history (O($N$) per step), which would compound the already significant throughput penalty. With group\_size=32, the amortized cost drops to O(1), making KIVI's overhead constant regardless of sequence length---an essential property for long-context deployment.

%% ============================================================
\section{Work in Progress}
%% ============================================================

\textbf{Additional methods.} Our midpoint covers quantization (KIVI) and dynamic selection (TopK). For the final report we will implement two additional methods: xKV~\cite{xkv} (per-head truncated SVD, representing low-rank compression) and SnapKV~\cite{snapkv} (attention-based one-shot token eviction). This will complete coverage across all four KV cache optimization families and enable cross-family comparisons---e.g., whether quantization and eviction degrade quality on different token types, or whether low-rank approximation and selection can be composed.

\textbf{CUDA kernel integration.} Our pure-PyTorch approach ensures fair algorithmic comparison but inflates throughput overhead (KIVI 76\% slower, TopK 11--24\% slower than what optimized kernels achieve). For the final report we plan to integrate the original Triton/CUDA kernels: KIVI's fused quantized attention (avoids dequantize-requantize overhead) and TokenSelect's paged sparse kernels. If framework incompatibilities prevent integration, our fallback is to benchmark each method's official open-source implementation independently while keeping model weights, prompts, seeds, and hardware identical---sacrificing the single-codebase property but preserving the controlled comparison.

\textbf{Longer and more realistic sequences.} Current synthetic prompts are WikiText passages truncated to controlled lengths. We plan to extend to 16K--32K tokens using naturally occurring long documents, increase the generation budget beyond 200 tokens for summarization, and evaluate with real-world attention patterns rather than artificial truncations.

\textbf{Verification pass.} The midpoint focus has been on understanding each method's algorithmic properties and obtaining an initial picture of the tradeoff space. Before the final submission we will audit implementations against original pseudocode and reference code, cross-check KV cache size accounting against actual GPU memory deltas, and re-run a subset of experiments to confirm reproducibility.

\textbf{Additional dimensions.} We plan to add kernel-level profiling via \texttt{torch.profiler} with TensorBoard trace export to identify whether KIVI's bottleneck lies in quantization arithmetic, memory allocation, or synchronization. We will also run batch size scaling experiments (1/4/8/16) to study how cache compression interacts with concurrent requests---KIVI should benefit disproportionately since per-request memory savings compound with batch size. Finally, we aim to evaluate hybrid method combinations (e.g., KIVI quantization applied to TopK's full cache) to test whether methods from different optimization families exhibit complementary benefits.

%% ============================================================
\section{Blockers and Limitations}
%% ============================================================

\textbf{Single model architecture.} All results are specific to LLaMA-2-7B with multi-head attention (32 heads, dim 128). Models using grouped-query attention (e.g., LLaMA-3, Mistral) have fewer KV heads, which would change compression characteristics. We plan to evaluate on at least one GQA model if time permits.

\textbf{Greedy decoding.} We use deterministic argmax decoding for reproducibility. Sampling-based generation (temperature, top-p) could interact differently with approximate KV caches: perturbations that leave the argmax unchanged may alter sampling outcomes, potentially amplifying quality differences between methods.

\textbf{LongBench score magnitudes.} Scores are modest across all methods including baseline (avg $\sim$0.06), reflecting the difficulty of these tasks for a 7B model with greedy decoding and a 200-token generation limit. The scores remain useful for relative comparison between methods, which is the intended purpose. For the final report, increasing the generation budget and evaluating with a larger model could yield more discriminative task-level scores.

\textbf{Pure-PyTorch throughput.} All implementations use standard PyTorch operations without custom CUDA or Triton kernels. This ensures fair algorithmic comparison---every method pays the same Python overhead---but means that absolute throughput numbers (especially KIVI's 76\% penalty) are not representative of production performance. The relative ordering by algorithmic complexity is expected to hold under kernel optimization, but the absolute gaps will narrow substantially.

\bibliographystyle{IEEEtran}

\begin{thebibliography}{10}

\bibitem{kivi}
Z.~Liu \textit{et al.}, ``KIVI: A tuning-free asymmetric 2bit quantization for KV cache,'' in \textit{ICML}, 2024.

\bibitem{topk}
Z.~Wu \textit{et al.}, ``TokenSelect: Efficient long-context inference via dynamic token-level KV cache selection,'' in \textit{EMNLP}, 2025.

\bibitem{xkv}
C.~Chang \textit{et al.}, ``xKV: Cross-layer SVD for KV-cache compression,'' \textit{arXiv:2503.18893}, 2025.

\bibitem{llama2}
H.~Touvron \textit{et al.}, ``Llama 2: Open foundation and fine-tuned chat models,'' \textit{arXiv:2307.09288}, 2023.

\bibitem{wikitext}
S.~Merity \textit{et al.}, ``Pointer sentinel mixture models,'' in \textit{ICLR}, 2017.

\bibitem{longbench}
Y.~Bai \textit{et al.}, ``LongBench: A bilingual, multitask benchmark for long context understanding,'' \textit{arXiv:2308.14508}, 2023.

\bibitem{snapkv}
Y.~Li \textit{et al.}, ``SnapKV: LLM knows what you are looking for before generation,'' in \textit{NeurIPS}, 2024.

\end{thebibliography}

\end{document}
